\documentclass[10pt,a4paper]{amsart}
\usepackage{amsmath,amssymb,amsthm,graphicx}
% Retained minimal preamble for the Stacks Project chapter more-algebra.tex (More on Algebra).
% The chapter file pinned-original-source.tex carries only \input{preamble} on its first line, so no
% preamble text lives inside the chapter itself.  The Stacks Project uses one single shared preamble for
% every chapter; the snapshot stages it identically as sources/stacks-chow/preamble.tex and
% sources/stacks-etale-cohomology/preamble.tex.  The block below is copied mechanically, line for line,
% from sources/stacks-chow/preamble.tex lines 14-18 (the Xy-pic setup) and 154-217 (the native theorem
% declarations together with the mathematical macros).  Nothing was added, removed or redefined; the
% wrapper class and the amsmath/amssymb/amsthm/graphicx base packages are supplied by the assembler.
% Chapter: more-algebra.tex -- More on Algebra.
% For commutative diagrams we use Xy-pic
\usepackage[all]{xy}

% We use 2cell for 2-commutative diagrams.
\xyoption{2cell}

% Theorem environments.
%
\theoremstyle{plain}
\newtheorem{theorem}[subsection]{Theorem}
\newtheorem{proposition}[subsection]{Proposition}
\newtheorem{lemma}[subsection]{Lemma}

\theoremstyle{definition}
\newtheorem{definition}[subsection]{Definition}
\newtheorem{example}[subsection]{Example}
\newtheorem{exercise}[subsection]{Exercise}
\newtheorem{situation}[subsection]{Situation}

\theoremstyle{remark}
\newtheorem{remark}[subsection]{Remark}
\newtheorem{remarks}[subsection]{Remarks}

\numberwithin{equation}{subsection}

% Macros
%
\def\lim{\mathop{\mathrm{lim}}\nolimits}
\def\colim{\mathop{\mathrm{colim}}\nolimits}
\def\Spec{\mathop{\mathrm{Spec}}}
\def\Hom{\mathop{\mathrm{Hom}}\nolimits}
\def\Ext{\mathop{\mathrm{Ext}}\nolimits}
\def\SheafHom{\mathop{\mathcal{H}\!\mathit{om}}\nolimits}
\def\SheafExt{\mathop{\mathcal{E}\!\mathit{xt}}\nolimits}
\def\Sch{\mathit{Sch}}
\def\Mor{\mathop{\mathrm{Mor}}\nolimits}
\def\Ob{\mathop{\mathrm{Ob}}\nolimits}
\def\Sh{\mathop{\mathit{Sh}}\nolimits}
\def\NL{\mathop{N\!L}\nolimits}
\def\CH{\mathop{\mathrm{CH}}\nolimits}
\def\proetale{{pro\text{-}\acute{e}tale}}
\def\etale{{\acute{e}tale}}
\def\QCoh{\mathit{QCoh}}
\def\Ker{\mathop{\mathrm{Ker}}}
\def\Im{\mathop{\mathrm{Im}}}
\def\Coker{\mathop{\mathrm{Coker}}}
\def\Coim{\mathop{\mathrm{Coim}}}

% Boxtimes
%
\DeclareMathSymbol{\boxtimes}{\mathbin}{AMSa}{"02}

%
% Macros for moduli stacks/spaces
%
\def\QCohstack{\mathcal{QC}\!\mathit{oh}}
\def\Cohstack{\mathcal{C}\!\mathit{oh}}
\def\Spacesstack{\mathcal{S}\!\mathit{paces}}
\def\Quotfunctor{\mathrm{Quot}}
\def\Hilbfunctor{\mathrm{Hilb}}
\def\Curvesstack{\mathcal{C}\!\mathit{urves}}
\def\Polarizedstack{\mathcal{P}\!\mathit{olarized}}
\def\Complexesstack{\mathcal{C}\!\mathit{omplexes}}
% \Pic is the operator that assigns to X its picard group, usage \Pic(X)
% \Picardstack_{X/B} denotes the Picard stack of X over B
% \Picardfunctor_{X/B} denotes the Picard functor of X over B
\def\Pic{\mathop{\mathrm{Pic}}\nolimits}
\def\Picardstack{\mathcal{P}\!\mathit{ic}}
\def\Picardfunctor{\mathrm{Pic}}
\def\Deformationcategory{\mathcal{D}\!\mathit{ef}}
\usepackage{cleveref}
\crefname{originallemma}{original lemma}{original lemmas}
\Crefname{originallemma}{Original lemma}{Original lemmas}
\makeatletter
\newlabel{more-algebra-lemma-reflexive-torsion-free}{{0AV0}{source}{}{sourcepointer-1}{}}
\newlabel{more-algebra-lemma-reflexive-torsion-free@cref}{{[originallemma][0][]0AV0}{[1][0][]source}{}{}{}}
\makeatother
\crefname{originallemma}{original lemma}{original lemmas}
\Crefname{originallemma}{Original lemma}{Original lemmas}
\makeatletter
\newlabel{lemma-hom-into-torsion-free}{{0AUX}{source}{}{sourcepointer-2}{}}
\newlabel{lemma-hom-into-torsion-free@cref}{{[originallemma][0][]0AUX}{[1][0][]source}{}{}{}}
\makeatother
\title{Reflexive modules are torsion free and the double dual map is injective exactly for finite torsion free modules}
\author{The Stacks Project Authors}
\date{stacks-project:more-algebra.tex:local-snapshot-20261012; GNU FDL 1.2 or later}
\begin{document}
\maketitle
\setcounter{subsection}{1}
\noindent\textit{Editorial note.} All delivered mathematical text is the author's own, copied from the pinned Stacks Project chapter file more-algebra.tex; no statement or proof step was rewritten, reordered or removed.\par\medskip
\noindent\textit{Editorial note.} The chapter file pinned-original-source.tex carries only \textbackslash{}input\{preamble\} on its first line, so no preamble text lives inside the chapter. The retained file sources/stacks-more-algebra/retained-preamble.tex is a mechanical, line-for-line copy of the project's single shared preamble as staged in the snapshot, namely sources/stacks-chow/preamble.tex lines 14-18 (Xy-pic setup) and 154-217 (theorem declarations and mathematical macros). Nothing was added, removed or redefined, no package was added, and the excerpt serves only as preamble, never as a span coordinate.\par\medskip
\noindent\textit{Editorial note.} Both counted spans start at the \textbackslash{}begin line of their environment, so the statement and the author's proof are delivered whole; surrounding prose of the section is deliberately outside the counted spans and is carried as standard prerequisites instead.\par\medskip
\noindent\textit{Editorial note.} The citation fields (source\_id, source\_version, authors, source\_url, license.name, source\_title) are identical to those of the cards already published from this chapter, so that the batch selector treats them as one source instead of deferring them as a same-source clash.\par\medskip
\noindent\textit{Allowed standard prerequisites (editorial).}\begin{itemize}
\item The double dual map j : M to Hom(Hom(M, R), R) and the definition of a reflexive module. definition environment at source lines 5148-5157, immediately before the counted statement. Supplies the map whose kernel, cokernel and injectivity the three assertions of the statement are about.
\item The module of maps into a torsion free module is torsion free, applied to the double dual of R. the author cites this at proof lines 5178-5179 as Lemma lemma-hom-into-torsion-free (tag 0AUX), the preceding lemma of the same chapter. Gives assertion (1) immediately and also shows injectivity of j for torsion free M in the last line of the proof.
\item Finitely generated modules over a domain have finite rank, and M tensor K is a finite dimensional vector space. used at proof lines 5180-5187 to choose the generators x\_1 to x\_r that map to a basis and to clear denominators. Produces the partial generating system and the element c with which the two comparison maps alpha and beta are built.
\end{itemize}\medskip
\ifdefined\hypertarget\hypertarget{sourcepointer-1}{}\fi
\noindent\textit{Source reference.} Original lemma 0AV0: focused-seeds/AGP/sources/stacks-snapshot-20261011/tags.txt, source lines 13668--13668. The lemma's own canonical tag, quoted so that readers can resolve the label inside the delivered text; it points at the counted statement itself and carries no proof content.\par\medskip
\ifdefined\hypertarget\hypertarget{sourcepointer-2}{}\fi
\noindent\textit{Source reference.} Original lemma 0AUX: focused-seeds/AGP/sources/stacks-snapshot-20261011/tags.txt, source lines 13665--13665. Navigation pointer to the immediately preceding lemma of the same chapter, which the author invokes for the torsion freeness of the double dual; that statement is counted by the sibling card S1-algebra-hom-into-torsion-free of this lane, and the construction with alpha, beta and the diagram remains inside the counted proof span.\par\medskip
\begin{lemma}
\label{lemma-reflexive-torsion-free}
Let $R$ be a domain and let $M$ be an $R$-module.
\begin{enumerate}
\item If $M$ is reflexive, then $M$ is torsion free.
\item If $M$ is finite, then the kernel and cokernel of
$j : M \to \Hom_R(\Hom_R(M, R), R)$ are torsion modules.
\item If $M$ is finite, then $j$ is injective
if and only if $M$ is torsion free.
\end{enumerate}
\end{lemma}

\begin{proof}
To see (1) use that the double dual of $M$ is torsion free by
Lemma \ref{lemma-hom-into-torsion-free}. Assume that $M$ is finite.
Choose generators $x_1, \ldots, x_n$ of $M$ over $R$.
Let $K$ be the fraction field of $R$. After renumbering we
may assume that for some $0 \leq r \leq n$ the elements
$x_1, \ldots, x_r$ map to a basis of $M \otimes_R K$ over $K$.
For $r < i \leq n$ write
$x_i \otimes 1 = \sum_{1 \leq j \leq r} x_j \otimes k_{ij}$ with
$k_{ij} \in K$. Choose $c \in R$ such that $a_{ij} = ck_{ij} \in R$
and such that $cx_i =  \sum_{1 \leq j \leq r} a_{ij} x_j$ in $M$
(small detail omitted). Then we get maps
$$
R^{\oplus r} \xrightarrow{\alpha} M \xrightarrow{\beta} R^{\oplus r}
\xrightarrow{\alpha} M
$$
with both $\alpha \circ \beta$ and $\beta \circ \alpha$
given by multiplication with $c$. Here $\alpha$ uses
$x_1, \ldots, x_r$ and $\beta$ sends $x_i$ to $c$ times the $i$th
basis vector for $1 \leq i \leq r$ and
to the vector $(a_{i1}, \ldots, a_{ir})$ for $r < i \leq n$.
Consider the commutative diagram
$$
\xymatrix{
R^{\oplus r} \ar[r] \ar[d]^{\text{id}} &
M \ar[r] \ar[d]^j &
R^{\oplus r} \ar[d]^{\text{id}} \ar[r] &
M \ar[d]^j \\
R^{\oplus r} \ar[r] &
\Hom_R(\Hom_R(M, R), R) \ar[r] &
R^{\oplus r} \ar[r] &
\Hom_R(\Hom_R(M, R), R)
}
$$
Of course the compositions of consecutive arrows on the bottom row
are equal to multiplication by $c$. The diagram shows that the kernel
and cokernel of $j$ are annihilated by $c$ which proves (2).
For (3) note that if $M$ is torsion free, then
$\beta$ is injective which implies that $j$ is injective.
The converse is clear because we've seen already that the double
dual of $M$ is torsion free.
\end{proof}

\end{document}
